% casting_notes.tex — Casting philosophy and notes

\section{Casting Notes}

\subsection{The Centerpiece: Adrien Brody as Kafka}

The script is engineered for Brody's specific instrument. The role demands:

\begin{itemize}[leftmargin=*]
\item Extended passages of minimal dialogue where the performance lives entirely in the face and body
\item Visible physical deterioration across the film's timeline (age 29 to 40, healthy to terminal)
\item The capacity to make shame, humor, bureaucratic precision, and intellectual exactness coexist within a single expression
\item A quality of inwardness that holds the camera without external action
\item Fluency in multiple emotional registers: dry comedy, social terror, romantic tenderness, philosophical combat, physical agony, and spiritual clarity
\item The ability to embody the role's central contradiction: a man who can make institutions tell the truth and cannot make his own father read a letter
\end{itemize}

Brody's combination of physical distinctiveness, emotional transparency, and demonstrated capacity for sustained character immersion makes this a natural fit at the highest level. Post-\textit{Brutalist}, the timing is precise.

\subsection{Supporting Cast Philosophy}

The supporting cast must be European or convincingly European. This is not a film that can absorb anachronistic casting. Period credibility and linguistic plausibility are essential.

\subsubsection{Hermann Kafka}
Requires a commanding physical presence and the capacity to be simultaneously charming and crushing. The actor must be likable --- the role's power comes from the fact that Hermann is not wrong about everything, and his love for Franz is real, even if it is expressed in a language Franz cannot speak. Consider actors in the Brendan Gleeson, Stellan Skarsgård register: large men with enormous range.

\subsubsection{Max Brod}
The foil and the only genuine philosophical rival in the film. Must be warm, energetic, slightly pushy, and capable of the moral complexity of the manuscripts decision. A character actor of real depth who can make love, opportunism, conviction, and disobedience exist at once. Consider the Mark Rylance, Ben Whishaw register: actors who combine intelligence with warmth.

\subsubsection{Felice Bauer}
Must feel practical, grounded, real. Not a period decoration. She represents structure, timetable, adult life, and a future Kafka cannot enter. The actor must convey genuine attraction, genuine bewilderment, professional competence, and the dignity of a woman who has been let down by someone she tried to love honestly.

\subsubsection{Milena Jesenská}
The most intellectually electric presence in the film after Kafka himself. She is not simply another love interest but the shock of being read accurately. Must convey brilliance, directness, erotic intelligence, and the specific magnetism of a woman who sees through people. European casting essential --- this character must feel Czech, not generically European.

\subsubsection{Dora Diamant}
Warmth without sentimentality. Practicality without coldness. The actor must convey what makes Dora different from every other person in Kafka's life: she offers presence without negotiation, which is why the Berlin movement can briefly feel like a real life instead of another variation on failure. Her tenderness in the final act must be devastating without being theatrical.

\subsubsection{Julie Kafka}
Warmth and tiredness and a moral clarity that arrives without announcement. Julie has spent thirty years carrying things between Hermann and the world, and between Hermann and Franz, and she does not eat either --- the same symptom, different cause, unnoticed because she is the one who does the feeding. The actress must convey the specific weight of a woman who refuses to deliver a hundred-page letter not through cowardice but through understanding what it would destroy, and who later mails it back to Berlin because she can no longer bear to serve as the family's permanent vault. Her two choices with the Letter define the title engine.

\subsubsection{Ottla Kafka}
The youngest sister, the family rebel, and the only person in Kafka's life who tells him the truth to his face without first arranging the furniture. She defied Hermann about the farm --- in person, not by letter --- which gives her a moral authority her brother lacks. The actress must be young, warm, direct, and capable of delivering the film's most lacerating line (``You wrote a hundred pages because paper doesn't interrupt'') with love rather than cruelty.

\subsection{Ensemble Approach}

The film's minor characters --- shop girls, factory workers, sanatorium patients, doctors, nurses, clerks, operators --- should be cast with the same care as the principals. Every face in the film tells a story. The world Kafka moves through must feel populated by real people, not extras, because the office, the family, and the medical system all need to feel like actual human environments rather than symbolic wallpaper.
